\section{Methodology}
\label{methodology}

\subsection{Framework Overview}
\label{sec:method-overview}

Given a set of natural-language requirements $R=\{r_1,\ldots,r_n\}$, the goal is to construct a structured requirement representation $G=(V,E,A)$.
Nodes $V$ represent requirement elements, edges $E$ represent semantic relations, and attributes $A$ store textual evidence, type labels, confidence, and human revision status.
The framework contains three stages: multi-agent parsing, knowledge-guided structure modeling, and human-in-the-loop refinement.

\subsection{Multi-Agent Requirement Parsing Workflow}
\label{sec:method-workflow}

\subsubsection{Interviewer Agent}
\label{sec:method-interviewer}

The interviewer agent rewrites stakeholder-facing requirement text into candidate requirement statements and records unresolved assumptions.
Its output is a cleaned requirement list rather than a final model.

\subsubsection{Analyst Agent}
\label{sec:method-analyst}

The analyst agent decomposes each candidate requirement into semantic units, including functions, constraints, roles, data elements, and exception information.
It also links each extracted unit back to textual evidence.

\subsubsection{Structure Modeling Agent}
\label{sec:method-structure-agent}

The structure modeling agent maps decomposed semantic units into the requirement graph.
It normalizes element types, constructs relations, and prepares the representation for human inspection.

\subsection{Requirement Semantic Decomposition}
\label{sec:method-decomposition}

\subsubsection{Functional Requirements}
\label{sec:method-functional}

Functional requirements describe intended system behaviors, services, or responses to user actions.

\subsubsection{Constraints}
\label{sec:method-constraints}

Constraints capture non-behavioral conditions that restrict how a function should operate, such as quality expectations, policy limits, or operational conditions.

\subsubsection{Roles}
\label{sec:method-roles}

Roles identify stakeholders, users, external systems, or organizational actors involved in a requirement.

\subsubsection{Data Elements}
\label{sec:method-data}

Data elements represent objects, records, attributes, and information items manipulated or referenced by a requirement.

\subsubsection{Exception Information}
\label{sec:method-exceptions}

Exception information captures alternative conditions, error cases, missing inputs, or boundary situations stated in the requirement text.

\subsection{Knowledge-Guided Structure Modeling}
\label{sec:method-knowledge}

\subsubsection{Requirement Knowledge}
\label{sec:method-req-knowledge}

Requirement knowledge includes domain vocabulary, requirement categories, expected element types, and common relation types.

\subsubsection{Structure Templates}
\label{sec:method-templates}

Structure templates define the expected schema for decomposed requirement elements and graph relations.

\subsubsection{Semantic Patterns}
\label{sec:method-patterns}

Semantic patterns guide the interpretation of common requirement forms, such as actor-action-object statements and condition-trigger-response statements.

\subsubsection{Prompt Strategies}
\label{sec:method-prompts}

Prompt strategies constrain agent outputs to typed fields, evidence links, and graph construction rules.

\subsection{Requirement Graph Construction}
\label{sec:method-graph}

\subsubsection{Entity Extraction}
\label{sec:method-entity}

Entity extraction identifies requirement elements and assigns each element a type, source requirement, and evidence span.

\subsubsection{Relation Construction}
\label{sec:method-relation}

Relation construction connects entities through typed links such as performs, uses, constrains, triggers, and handles.

\subsubsection{Structured Representation}
\label{sec:method-structured-representation}

The final representation stores the requirement graph in a structured format that can be inspected, compared, and extended in later requirements engineering activities.

\subsection{Human-in-the-Loop Refinement}
\label{sec:method-hitl}

Human reviewers inspect uncertain nodes, ambiguous relations, and unsupported decompositions.
Their edits are recorded as revisions to the graph rather than as unstructured comments, keeping the representation analyzable after refinement.
